% Author: Dominik Harmim <xharmi00@stud.fit.vutbr.cz>


\documentclass[a4paper, 11pt]{article}


\usepackage[czech]{babel}
\usepackage[utf8]{inputenc}
\usepackage[left=2cm, top=3cm, text={17cm, 24cm}]{geometry}
\usepackage{times}
\usepackage{graphicx} % vkládání obrázků
\usepackage[unicode]{hyperref}
\hypersetup{
	colorlinks = true,
	hypertexnames = false,
	citecolor = red
}


\newcommand{\RNum}[1]{\uppercase\expandafter{\romannumeral #1\relax}} % makro na sázení římských čísel


\begin{document}

	%%%%%%%%%%%%%%%%%%%%%%%%%%%%%%%% Titulní stránka %%%%%%%%%%%%%%%%%%%%%%%%%%%%%%%%
	\begin{titlepage}
		\begin{center}
			\includegraphics[width=0.77\linewidth]{inc/FIT_logo.pdf} \\

			\vspace{\stretch{0.382}}

			\Huge{Projektová dokumentace} \\
			\LARGE{\textbf{Implementace překladače imperativního jazyka IFJ17}} \\
			\Large{Tým 104, varianta \RNum{2}}
			\vspace{\stretch{0.618}}
		\end{center}

		\begin{minipage}{0.4 \textwidth}
			{\Large \today}
		\end{minipage}
		\hfill
		\begin{minipage}[r]{0.6 \textwidth}
			\Large
			\begin{tabular}{l l l}
				\textbf{Dominik Harmim} & \textbf{(xharmi00)} & \quad 25\,\% \\
				Vojtěch Hertl & (xhertl04) & \quad 25\,\% \\
				Timotej Halás & (xhalas10) & \quad 25\,\% \\
				Matej Karas & (xkaras34) & \quad 25\,\% \\
			\end{tabular}
		\end{minipage}
	\end{titlepage}


	%%%%%%%%%%%%%%%%%%%%%%%%%%%%%%%% Obsah %%%%%%%%%%%%%%%%%%%%%%%%%%%%%%%%
	\pagenumbering{roman}
	\setcounter{page}{1}
	\tableofcontents
	\clearpage


	%%%%%%%%%%%%%%%%%%%%%%%%%%%%%%%% Úvod %%%%%%%%%%%%%%%%%%%%%%%%%%%%%%%%
	\pagenumbering{arabic}
	\setcounter{page}{1}

	\section{Úvod}

	Cílem projektu bylo vytvořit program v~jazyce~C, který načte zdrojový kód zapsaný ve zdrojovém jazyce IFJ17,
	jenž je zjednodušenou podmnožinou jazyka FreeBASIC a~přeloží jej do cílového jazyka IFJcode17 (mezikód).

	Program funguje jako konzolová aplikace, které načítá zdrojový program ze standardního vstupu a~generuje
	výsledný mezikód na standardní výstup nebo v~případě chyby vrací odpovídající chybový kód.


	%%%%%%%%%%%%%%%%%%%%%%%%%%%%%%%% Návrh a implementace %%%%%%%%%%%%%%%%%%%%%%%%%%%%%%%%
	\section{Návrh a~implementace}

	Projekt jsme sestavili z~několika námi implementovaných dílčích částí, které jsou představeny v~této kapitole.
	Je zde také uvedeno, jakým způsobem spolu jednotlivé dílčí části spolupracují.


	\subsection{Lexikální analýza}

	Při tvorbě překladače jsme začali implementací lexikální analýzy. Hlavní funkce této analýzy je \texttt{get\_next\_token}, 
	pomocí níž se čte znak po znaku ze zdrojového souboru a~převádí na strukturu token, která se skládá z~typu a~atributu. 
	Typy tokenu jsou \texttt{EOF}, \texttt{EOL}, prázdný token, identifikátory, klíčová slova, celé či desetinné číslo, řetězec 
	a~také porovnávací a~aritmetické operátory a~ostatní znaky, které mohou být použity v~jazyce IFJ2017. Typ atributu je 
	\texttt{union} a~dělí se podle typu tokenu. Pokud je typ tokenu řetězec nebo identifikátor, pak bude atribut daný řetězec, 
	když by byl typ tokenu klíčové slovo, přiřadí atributu dané klíčové slovo, pokud číslo, atribut bude ono číslo, u~ostatních 
	případů není atribut použit. S~takto vytvořeným tokenem poté pracují další analýzy.

	Celý lexikální analyzátor je implementován jako deterministický konečný automat podle předem vytvořené-ho diagramu 
	\ref{figure:fa_graph}. Konečný automat je v~jazyce C~jako jeden nekonečně opakující se \texttt{switch}, kde každý případ 
	\texttt{case} je ekvivalentní k~jednomu stavu automatu. Pokud načtený znak nesouhlasí s~žádným znakem, který jazyk povoluje, 
	program je ukončen a~vrací chybu 1. Jinak se přechází do dalších stavů a~načítají se další znaky, dokud nemáme hotový jeden 
	token, který potom vracíme a~ukončíme tuto funkci. 

	Za zmínku stojí zpracovávání identifikátorů a~klíčových slov, kde načítáme znaky a~průběžně ukládáme do dynamického pole 
	a~jakmile odcházíme ze stavu identifikátoru, porovnáváme, jestli načtený řetězec není shodný s~nějakým z~klíčových slov a~podle 
	toho se rozhodneme, zda vracíme atribut a~typ tokenu klíčové slovo či identifikátor. 
	
	Pro zpracovávání escape sekvencí v~řetězci máme vytvořeno pole o velikosti 4, které je zpočátku vynulováno a poté se postupně 
	naplňuje přečtenými čísly, vždy na pozici podle toho, kolikáté číslo zrovna uvažujeme a nakonec celé třímístné čislo převedeme 
	do ASCII podoby.


	\subsection{Syntaktická analýza}

	TODO

	\subsubsection{Syntaktická analýza založená na LL -- gramatice využívající metodu rekurzivního sestupu}

	TODO

	\subsubsection{Zpracování výrazů pomocí precedenční syntaktické analýzy}

	TODO


	\subsection{Sémantická analýza}

	TODO


	\subsection{Generování cílového kódu}

	Generování cílového kódu pro nás znamená generování mezikódu IFJcode17. Kód je generován na standardní výstup
	po dokončení všech analýz, přičemž jednotlivé dílčí části jsou generovány v~průběhu analýz do interní paměti
	programu. Při generování mezikódu jsou generovány i~komentáře pro jeho zpřehlednění.

	\subsubsection{Rozhraní generátoru kódu}

	Generování kódu je implementováno jako samostatný modul v~souboru \texttt{code\_generator.c},
	jehož rozhraní se nachází v~souboru \texttt{code\_generator.h} Rozhraní nabízí 3 základní funkce.
	Jedná se o~funkci, která připraví generování kódu, tato funkce se volá na začátku syntaktické analýzy.
	Funkce, která uvolní z~paměti všechny rezervované prostředky generátorem kódu se volá po dokončení
	syntaktické analýzy. Poslední důležitá funkce generuje vytvořený kód na patřičný výstup. Rozhraní
	nabízí spoustu dalších funkcí pro generování jednotlivých částí programu, např. generování začátku
	funkce, volání funkce, deklarace proměnnné, příkazu \texttt{input}, funkcí pro zpracování výrazů,
	podmínek, cyklů, aj. Všechny tyto funkce se volají ve správnou chvíli a~se správnými
	parametry v~průběhu syntaktické analýzy a~syntaktické analýzy výrazů.

	\subsubsection{Začátek generování}

	Na začátku generování jsou inicializovány potřebné datové struktury (které jsou na závěr
	opět uvolněny), vygenerována hlavička mezikódu, která zahrnuje potřebné náležitosti pro korektní
	interpretaci mezikódu, definici globálních proměnných na globálním rámci a~skok do hlavního
	těla programu. Poté jsou vygenerovány vestavěné funkce, které jsou zapsány přímo v~jazyce IFJcode17.

	\subsubsection{Generování funkcí}

	Každá funkce je tvořena návěštím ve tvaru \texttt{\$identifikator\_funkce} a~svým lokálním rámcem.
	Pro funkce se vždy generuje výchozí návratová hodnota podle jejího návratového datového typu,
	která se ukládá na lokální rámec, který je po odchodu z~funkce dostupný jako dočasný rámec.
	Před zavoláním funkce jsou hodnoty parametrů uloženy do dočasného rámce s~číslem, které
	odpovídá pořadí daného parametru a~po vstupu do funkce jsou všechny předané parametry uloženy na
	lokální rámce se svým názvev. Při generování návratu z~funkce je předán výsledek zpracovaného
	výrazu do proměnné s~návratovou hodnotou funkce a~proveden skok na konec funkce.

	\subsubsection{Generování výrazů}

	Všechny výrazy jsou během syntaktické analýzy výrazů ukládány na datový zásobník a~v~pravou
	chvíli jsou provedeny patřičné operace s~hodnotami na vrcholu zásobníku. Hodnota zpracovaného
	výrazu se ukládá do globální proměnné na globálním rámci.

	\subsubsection{Generování návěští}

	Všechna návěští, pokud neuvažujeme návěští pro funkce a~jiná speciální návěští, jsou generovány
	ve tvaru \texttt{\$identifikator\_funkce\%label\_deep\%label\_index}, kde
	\texttt{identifikator\_funkce} zajistí unikátnost návěští mezi jednotlivými funkcemi,
	\texttt{label\_deep} je hloubka zanoření návěští (například u~podmínek a~cyklů)
	a~\texttt{label\_index} je číselná hodnota, která se inkrementuje s~každým dalším návěštím uvnitř
	dané funkce.


	\subsection{Překladový systém}

	Projekt je možné přeložit dvěma způsoby, buď nástrojem CMake nebo nástrojem GNU Make. Oba způsoby vytvoří
	spustitelný soubor \texttt{ifj17\_compiler}.

	\subsubsection{CMake}

	Při vývoji, ladění a~testování projketu někteří z~nás používali CMake, protože překlad tímto nástrojem
	je integrován ve vývojovém prostředí CLion, které většina členů týmu používala a~protože nastavení pravidel
	pro překlad tímto nástrojem je jednoduché.

	V~souboru \texttt{CMakeLists.txt} jsou nastavena pravidla pro překlad. Je zde nastaven překladač \texttt{gcc}
	a~všechny potřebné parametry pro překlad. Dále je zde uvedeno, že se pro překlad mají používat všechny soubory s~příponou
	\texttt{.c} a~\texttt{.h} a~taky je zde uveden název spustitelného souboru.

	Soubory pro překlad je možné vygenerovat příkazem \texttt{cmake~.} a~překlad provést příkazem \texttt{cmake~--build~.}
	nebo pro toto využít nástroje vývojového prostředí.

	\subsubsection{GNU Make}

	Jedním z~požadavků pro odevzdání projektu bylo přiložit k~odevzdanému projektu i~soubor \texttt{Makefile} a~přeložení
	projektu příkazem \texttt{make}. Z~tohoto důvodu jsme vytvořili překladový systém i~pomocí nástroje GNU Make.

	V~souboru \texttt{Makefile} jsou nastavena pravidla pro překlad. Je zde nastaven překladač \texttt{gcc}
	a~všechny potřebné parametry pro překlad. Vytvořili jsme pravidla, které nejdříve vytvoří objektové soubory
	ze všech souborů s~příponou \texttt{.c} v~daném adresáři za pomoci vytvořených automatických pravidel a~překladače \texttt{gcc},
	který dokáže automaticky generovat tyto pravidla, více viz \emph{Dependecy Generation} v~kapitole \emph{C and C++}
	knihy \cite{Mecklenburgc2005}. Následně je ze všech objektových souborů vytvořen jeden spustitelný
	soubor.

	Nástroj GNU Make jsme mimo překlad využili i~pro automatické zabalení projektu pro odezvdání, generování dokumentace,
	spouštění automatických testů a~čištění dočasných souborů.


	%%%%%%%%%%%%%%%%%%%%%%%%%%%%%%%% Speciální algoritmy a datové struktury %%%%%%%%%%%%%%%%%%%%%%%%%%%%%%
	\section{Speciální algoritmy a~datové struktury}

	Při tvorbě překladače jsme implementovali několik speciálních datových struktur,
	které jsou v~této kapitole představeny.


	\subsection{Tabulka s~rozptýlenými položkami}

	Implementovali jsme tabulku s~rozptýlenými položkami, která slouží jako tabulka symbolů, což byl mimo
	jiné i~požadavek vyplývající ze zadání, který souvisí s~předmětem IAL. Tabulku jsme implementovali
	jako tabulku s~explicitním zřetězením synonym, protože jsome potřebovali, aby počet položek,
	který do ní můžeme uložit byl dynamický a~teoreticky neomezený. Pro zřetězení synonym jsme použili
	lineárně vázané seznamy. Velikost mapovacího pole jsme dimenzovali tak, aby byla rovna prvočíslu a~její
	naplnění nepřesáhlo 75\,\%. Jako optimální velikost jsme tedy zvolili číslo 27\,457.

	Jako mapovací funkci jsme použili GNU ELF Hash použitou v~algoritmu, který používá UNIX ELF.
	Tato funkce je variantou PJW Hash. Použití této mapovací funkce jsme vyhodnotili jako nejvhodnější,
	více viz \cite{PJWHashFunction2017}. Funkci jsme upravili tak, aby vracela datový typ a~hodnotu,
	které je přípustná pro naše mapovací pole.

	Každá položka tabulky obsahuje unikátní klíč v~podobě řetězce, kterým je identifikátor funkce
	či proměnné. Dále obsahuje ukazatel na další prvek (synonymum) a~datovou část. Datová část obsahuje
	informaci o~tom, jakého datového typu je daný symbol (v~případě funkce se jedná o~datový typ
	návratové hodnoty), pravdivostní hodnoty, které indikují, jestli byl daný symbol již definován
	a~zda je symbol globální. Posledním prvkem datové části, který se používá pouze v~případě, že se
	jedná o~funkci, je počet a~datové typy formálních parametrů funkce, který je ve formátu řetezce
	(Pokud např. daná funkce má dva parametry, první typu \texttt{Double} a~druhý typu \texttt{Integer},
	uloží se následující řetězec \texttt{di}.).

	Implementovali jsme i~několik potřebných funkcí, které tvoří rozhraní pro práci s~tabulkou.
	Jedná se o~následující funkce: inicializace, přidání nové položky, přidání určitého parametru
	pro konkrétní položku tabulky (v případě, že se jedná o~funkci), vyhledání položky, odstranění
	položky a~uvolnění tabulky z paměti. Při implementaci jsme se inspirovali studijní literaturou
	předmětu IAL \cite{Honzik1991}.


	\subsection{Dynamický řetězec}

	Vytvořili jsme strukturu \texttt{struct dynamic\_string} pro práci s~řetězci dynamické délky, kterou používáme
	např. pro ukládání řezězce nebo identifikátoru u~atributu tokenu v~lexikální analýze, protože dopředu nevíme,
	jak bude tento řetězec dlouhý.

	Daná struktura a~rozhraní operací nad ní je popsáno v~souboru \texttt{dynamic\_string.h} a~operace jsou implementovány
	v~souboru \texttt{dynamic\_string.c}.

	Struktura v~sobě uchovává ukazatel na řezězec, velikost řetězce a~informaci o~tom, kolik paměti je pro
	řetězec vyhrazeno. Implementovali jsme operace pro inicializaci struktury, uvolnění vyhrazených dat,
	přidání znaku na konec řezězce, porovnávání řetězců, kopie celých řezězců a~další. Při inicializaci struktury
	se vyhradí paměťový prostor pouze pro určitý počet znaků (určeno konstantou), až při nedostatku vyhrazené
	paměti se velikost paměti zvýší.


	\subsection{Zásobník symbolů pro precedenční syntaktickou analýzu}

	TODO


	%%%%%%%%%%%%%%%%%%%%%%%%%%%%%%%% Práce v týmu %%%%%%%%%%%%%%%%%%%%%%%%%%%%%%%%
	\section{Práce v~týmu}

	\subsection{Způsob práce v~týmu}

	Na projektu jsme začali pracovat koncem října. Práci jsme si dělili postupně, tj. neměli jsme od začátku
	stanovený kompletní plán rozdělení práce. Na dílčích částech projektu pracovali většinou jednotlivci nebo
	dvojice členů týmu. Nejprve jsme si stanovili strukturu projektu a~vytvořili překladový systém.

	\subsubsection{Verzovací systém}

	Pro správu souborů projektu jsme používali verzovací systém Git. Jako vzdálený repositář jsme používali \mbox{GitHub}.

	Git nám umožnil pracovat na více úkolech na projektu současně v~tzv. větvích. Většinu úkolů jsme nejdříve připravili
	do větve a~až po otestování a~schválení úprav ostatními členy týmu jsme tyto úpravy začlenili do hlavní
	vývojové větve.

	\subsubsection{Komunikace}

	Komunikace mezi členy týmů probíhala převážně osobně nebo prostřednictvím aplikace Slack, kde jsme si buď
	psali přímo mezi sebou nebo si vytvářeli různé skupinové konverzace podle tématu.

	V~průběhu řešení projektu jsme měli i~několik osobních setkání, kde jsem probírali a~řešili problémy
	týkající se různých částí projektu.


	\subsection{Rozdělení práce mezi členy týmu}

	Práci na projektu jsme si rozdělili rovnoměrně s~ohledem na její složitost a~časovou náročnost.
	Každý tedy dostal procentuální hodnocení 25\,\%.
	Tabulka \ref{table:rozdeleni_prace} shrnuje rozdělení práce v~týmu mezi jednotlivými členy.
	\bigskip
	\begin{table}[ht]
		\centering
		\begin{tabular}{| l | l |}
			\hline
			Člen týmu & Přidělená práce \\ \hline
			\textbf{Dominik Harmim} & \begin{tabular}{l} vedení týmu, organizace práce, dohlížení na provádění práce,
				konzultace, \\ kontrola, testování, dokumentace, struktura projektu, generování cílového kódu \end{tabular} \\
			Vojtěch Hertl & \begin{tabular}{l} lexikální analýza, testování, dokumentace, prezentace \end{tabular} \\
			Timotej Halás & \begin{tabular}{l} syntaktická analýza, syntaktická analýza výrazů, sémantická analýza, testování, \\
				 dokumentace \end{tabular} \\
			Matej Karas & \begin{tabular}{l} syntaktická analýza, sémantická analýza, testování, dokumentace \end{tabular} \\ \hline
		\end{tabular}
		\caption{Rozdělení práce v~týmu mezi jednotlivými členy}
		\label{table:rozdeleni_prace}
	\end{table}


	%%%%%%%%%%%%%%%%%%%%%%%%%%%%%%%% Závěr %%%%%%%%%%%%%%%%%%%%%%%%%%%%%%%%
	\section{Závěr}

	Projekt nás zprvu trochu zaskočil svým rozsahem a~složitostí. Postupem času, až jsme získali dostatek
	znalostí o~tvorbě překladačů na přednáškách IFJ, jsme projekt začali řešit.

	Náš tým jsme měli sestaven velmi brzy, byli jsme již předem domluveni na komunikačních kanálech,
	osobních schůzkách a~na používání verzovacího systému, tudíž jsme s~týmovou prací neměli žádný
	problém a~pracovalo se nám společně velmi dobře.

	Na projektu jsme začali pracovat trochu později, takže jsme neměli časovou rezervu, ale nakonec
	jsme všechno bez problému stihli. Jednotlivé části projektu jsme řešili většinou individuálně za použití
	znalostí z~přednáškek nebo materiálů do předmětů IFJ a~IAL.

	V~průběhu vývoje jsme se potýkali s~menšími problémy týkajícími se nejasností v~zadání, ale
	tyto jsme vyřešili díky fóru k~projektu. Správnost řešení jsme si ověřili automatickými
	testy a~pokusným odevzdáním, díky čemuž jsme byli schopni projekt ještě více odladit.

	Tento projekt nám celkově přinesl spoustu znalostí ohledně fungování překladačů, prakticky nám
	objasnil probíranou látku v~předmětech IFJ a~IAL a~přinesl nám zkušennosti s~projekty tohoto rozsahu.



	%%%%%%%%%%%%%%%%%%%%%%%%%%%%%%%% Citace %%%%%%%%%%%%%%%%%%%%%%%%%%%%%%%%
	\clearpage
	\bibliographystyle{czechiso}
	\renewcommand{\refname}{Literatura}
	\bibliography{dokumentace}


	%%%%%%%%%%%%%%%%%%%%%%%%%%%%%%%% Přílohy %%%%%%%%%%%%%%%%%%%%%%%%%%%%%%%%
	\clearpage
	\appendix

	\section{Diagram konečného automatu specifikující lexikální analyzátor}

	\begin{figure}[!ht]
		\centering
		\vspace{-1.2cm}
		\includegraphics[width=0.95\linewidth]{inc/FA_graph.pdf}
		\caption{Diagram konečného automatu specifikující lexikální analyzátor}
		\label{figure:fa_graph}
	\end{figure}
	\clearpage

	\section{LL -- gramatika}

	TODO


	\section{LL -- tabulka}

	TODO


	\section{Precedenční tabulka}

	TODO

\end{document}
